\section{Discussion}
\label{sec:discussion}

MouseBrainBench addresses a narrow but consequential problem: computational
neuroscience increasingly produces models that predict neural measurements,
while the language used to describe those models often extends to biological
mechanism or digital-twin fidelity. The experiments show that these
interpretations can be separated operationally. Known-truth systems expose
which evidential component is absent, Allen preserves a strong reproducibility
result while rejecting topology and direction, Sensorium preserves predictive
utility while rejecting causal interpretation, and MICRONS admits a local
structure--function association after matched controls and replication.

\subsection{Comparison with existing benchmark and validation paradigms}
\label{subsec:benchmark-comparison}

MouseBrainBench differs from predictive neuroscience benchmarks primarily in
its unit of evaluation. Sensorium evaluates a model on a defined prediction
task and provides a common basis for comparing architectures. That design is
appropriate for determining which predictor generalizes better under the
competition protocol \citep{willeke2022sensorium,turishcheva2023dynamic}.
MouseBrainBench instead evaluates the relationship between an empirical result
and the interpretation attached to it. Predictive correlation is retained as
one observation, but the benchmark instance also records whether repeated
responses are reliable, whether a proposed topology outperforms plausible
alternatives, and whether directional or interventional information is
observable. Consequently, the framework does not replace a leaderboard. It
determines which conclusions may legitimately be exported from one.

The distinction from simulation infrastructures is equally important. BMTK,
SONATA-compatible workflows, The Virtual Brain, and The Virtual Mouse Brain
provide mature mechanisms for constructing networks, executing dynamics, and
exchanging model descriptions \citep{dai2020bmtk,tvb2013,melozzi2017virtualmouse}.
Their central engineering object is an executable model. The central object in
MouseBrainBench is an evidence contract. A simulator can therefore be connected
as a candidate model without reproducing its numerical engine, but successful
execution is not itself treated as biological validation. Model outputs must
still be compared with observables, controls, and independent data appropriate
to the proposed claim.

This position is consistent with the broader verification, validation, and
uncertainty-quantification literature, which treats model credibility as
purpose-dependent rather than absolute \citep{riedmaier2021vvuq}. Digital-twin
credibility frameworks similarly emphasize intended use, lifecycle evidence,
data provenance, and uncertainty rather than assuming that the twin label
guarantees fidelity \citep{shao2022credibility}. MouseBrainBench contributes a
domain-specific operationalization for AI-based neural models: it maps purpose
to required evidence blocks, makes those blocks non-interchangeable, and binds
the reported interpretation to sealed computational artifacts. It does not yet
provide a complete uncertainty-quantification treatment of parameters and
model form. Instead, threshold sensitivity, matched nulls, clustered intervals,
and hold-out replication cover the uncertainties directly estimable from the
present resources.

Four practical properties therefore delimit the methodological contribution.
First, the benchmark compares claims rather than forcing heterogeneous models
onto one scalar leaderboard. Second, evidence types retain their semantics:
prediction, reproducibility, topology, direction, and structure--function
association cannot compensate for one another. Third, failed criteria remain
publishable outputs, making negative controls part of validation rather than
discarded development information. Fourth, the publication freeze joins
software revision, data identifiers, thresholds, and admissible language. No
single property is unprecedented in isolation. The contribution lies in their
executable combination across predictive, simulation, and connectomic regimes.

The negative Allen result is scientifically useful. A less restrictive
benchmark might report cross-mouse correlation of 0.891 and split-half
correlation of 0.990 as validation of a model. MouseBrainBench instead asks
whether the Allen-informed topology performs better than alternatives and
whether a directed signature can be recovered. It does not. This outcome does
not contradict the biological organization of the visual system. It identifies
an observational limitation: repeatable population dynamics do not uniquely
determine the graph that generated them.

Sensorium illustrates a related distinction in modern Artificial Intelligence.
Static stimulus and context models improve prediction over mean and scrambled
controls, while temporal models and a compact neural network improve predictions
for Dynamic Sensorium. These are valid functional results. They are not evidence
that the architecture's internal variables correspond to synapses, cell types,
or causal operations in the animal. Treating predictive and mechanistic scores
as separate outputs allows the same model to be accepted for one purpose and
rejected for another without contradiction.

The bounded official Sensorium experiment also clarifies what software
integration demonstrates. Executing an official loader and architecture verifies
interoperability, but a restricted local training budget is not comparable with
a published competition result. The very low bounded correlations are therefore
neither evidence against the official architecture nor a competitive baseline.
They document that the integration path works and that a full-budget comparison
remains absent. This prevents an engineering smoke test from becoming an
inflated empirical claim.

The MICRONS result is the strongest positive outcome, but its novelty must be
stated precisely. Ding et al. already demonstrated like-to-like connectivity
with a more extensive biological analysis, including feature, receptive-field,
and axon--dendrite geometry controls. MouseBrainBench does not improve that
biological finding. Instead, the MICRONS experiment tests whether a compact,
predefined audit can recover a compatible local association and preserve its
proper scope. Replication across two non-overlapping windows and positive
unit-cluster bootstrap intervals provide external validity for the audit
pipeline.

The effect is not causal. Readout-location similarity can covary with cortical
area, layer, cell class, developmental organization, or unmeasured morphology.
Distance and degree matching address two major confounds but do not exhaust the
causal graph. In particular, the analysis does not compare axon--dendrite
opportunity at the resolution used by Ding et al., manipulate individual
connections, or demonstrate that the observed edges are necessary for the
functional descriptor. The permitted interpretation is therefore a replicated
local observational association.

\subsection{Threats to validity and transfer beyond the mouse brain}
\label{subsec:validity-transfer}

The first threat concerns construct validity. MIS operationalizes mechanistic
identifiability through evidence blocks selected for the claim, but these blocks
are not a universal definition of mechanism. Latency and lead--lag information
are meaningful for the regional temporal systems examined here, whereas a
molecular pathway, a biophysical cell model, or a behavioural controller would
require different observables and alternatives. A passing decision means that
the declared operational criteria have been satisfied. It does not establish
that every scientifically plausible mechanism has been excluded. This is why
the continuous MIS value is diagnostic and why the categorical interpretation
must always be read together with the contract.

Internal validity is limited by the controls that can be constructed from each
resource. Permutations and disconnected models remove selected organizational
properties but cannot represent every latent circuit compatible with Allen
population activity. In MICRONS, distance and degree matching substantially
improve on random-pair comparison, yet cell class, laminar position, cortical
area, morphology, developmental lineage, and axon--dendrite opportunity may
still confound functional similarity. The aggregate analysis demonstrates the
importance of this issue: a positive random-pair contrast disappeared after
distance matching. The narrower replicated endpoint reduces the overclaim but
does not turn an observational association into a causal effect.

Statistical conclusion validity is constrained by the number and independence
of biological units. Pairwise MICRONS observations share neurons, so the study
uses unit-cluster rather than naive pair bootstrap intervals. Nevertheless, the
three cohorts are disjoint windows from one animal, not independent biological
replicates. Dynamic Sensorium comparisons contain five mice per cohort. Win
counts and medians are consequently more informative than asymptotic claims.
The sensitivity analysis reduces dependence on one threshold, but its tested
grid is finite and does not replace external preregistration. Multiple
descriptor and stratum tests are false-discovery-rate corrected, while the main
claim is restricted to the fixed endpoint replicated in both hold-outs.

External validity is bounded by species, brain area, task, and measurement
modality. Allen and Sensorium emphasize the mouse visual system. MICRONS covers
a restricted visual-cortical volume from one mouse, and the Allen experiments
use controlled subsets of the available sessions. The framework has not been
validated on subcortical circuits, longitudinal disease trajectories, freely
moving behaviour, closed-loop perturbation, or individual calibration across
modalities. These omissions directly block whole-brain and individual-twin
claims. They also mean that the current thresholds should not be transferred
unchanged to a different species or recording technology.

What can transfer is the evaluation grammar. A domain outside mouse neuroscience
can define a target claim, list the quantities that make it observable, specify
non-compensatory evidence blocks, construct claim-relevant alternatives, and
seal the resulting decision with provenance. In human neuroimaging, for
example, subject-level calibration and test--retest reliability might replace
session-level criteria, while interventions may remain unavailable. In
manufacturing, sensor calibration, operating-envelope coverage, and uncertainty
propagation would become mandatory blocks. In climate or physiological models,
conservation laws and out-of-regime validation could replace synaptic topology.
The software interfaces can represent these contracts, but scientific validity
still requires domain experts to justify the observables, controls, thresholds,
and consequences of failure.

The current MIS should also not be treated as a universal scalar definition of
mechanistic explanation. Its principal contribution is structural: evidence
blocks are non-interchangeable and passage criteria are explicit. The particular
criteria and thresholds remain claim-dependent. A molecular mechanism, a
neural-mass model, and an individual-level behavioural twin would require
different blocks. External preregistration and application by independent groups
are needed before MIS can be regarded as a community standard.

Conjunctive rules have an intentional cost. They reduce false promotion of
partial evidence but can reject scientifically interesting models when one
required measurement is unavailable. MouseBrainBench handles this by retaining
block scores and partial-positive labels. Failure of the global decision means
that the target claim is unsupported under the declared evidence. It does not
mean that the model, data, or partial result is useless. This semantic discipline
is important if the framework is to encourage stronger experiments rather than
merely produce negative labels.

The sensitivity analysis indicates that the main qualitative conclusions are
not artifacts of one exact threshold. Allen remains negative under 25\%
relaxation and tightening. Static Sensorium retains partial reliability and
topographic evidence for six of nine threshold combinations, failing when the
reliability threshold reaches 0.70. Dynamic Sensorium gains over mean response
are more stable than gains over temporal SVD, supporting predictive utility but
not superiority of a specific neural architecture. These patterns make the
limitations visible instead of hiding them in a single aggregate score.

For an engineering contribution in artificial intelligence, the value of the
work depends on the framework being useful beyond this exact collection of
analyses. The software interfaces are therefore organized around claims,
evidence blocks, graphs, predictions, perturbations, and provenance rather than
around hard-coded dataset names. External BMTK, SONATA, or The Virtual Brain
simulations can be added as candidate models without reimplementing their
numerical engines. New datasets can add evidence blocks while retaining the
same publication-freeze contract.

The next scientific step is not to increase model complexity indiscriminately.
It is to add evidence unavailable here: independent-animal structure--function
replication, perturbations that distinguish competing mechanisms, and
full-budget predictive baselines. Those additions would test whether the audit
can move from observational association towards intervention-supported
mechanism. Until then, MouseBrainBench should be understood as a reproducible
claim-control framework for partial mouse-brain models, not as the mouse-brain
model itself.
